% 表紙と目次
\thispagestyle{cover}
\colorlet{phase}{mone}\colorlet{phasetint}{monetint}
\begin{tcolorbox}[enhanced,colback=mone,colframe=mone,arc=3mm,boxrule=0pt,
    left=8mm,right=8mm,top=7mm,bottom=7mm,fontupper=\color{white}\sffamily]
  {\small 名古屋市立大学　2026年度 後期　土曜 1・2限　オンライン}\par
  \vspace{2mm}{\Huge\bfseries 経営情報 2026}\par
  \vspace{2mm}{\LARGE\bfseries 学習ガイド}\par
  \vspace{2.5mm}{\normalsize 全15回の講義予定に沿った内容解説・演習課題・模範解答・出席確認クイズの解説}\par
\end{tcolorbox}

\vspace{5mm}
{\sffamily\par\noindent
\begin{tblr}{colspec={Q[l,wd=24mm] X[l]},rowsep=2pt}
  著者 & 河合勝彦 \\
  所属 & 名古屋市立大学大学院経済学研究科 \\
  連絡先 & \href{mailto:kkawai@econ.nagoya-cu.ac.jp}{kkawai@econ.nagoya-cu.ac.jp} \\
  講義サイト & \href{https://kklab.mobi/kklab-mis2026/}{https://kklab.mobi/kklab-mis2026/} \\
  版 & 2026年10月9日版 \\
\end{tblr}}

\vspace{4mm}
\begin{tcolorbox}[enhanced,colback=monetint,colframe=monetint,arc=1.5mm,boxrule=0pt,
    left=4mm,right=4mm,top=2mm,bottom=2mm,fontupper=\small]
  {\sffamily\bfseries\color{mone}このガイドの構成}\par
  各回は「この回のねらい」「解説」「演習課題」「模範解答」「出席確認クイズの解説」の順に並んでいます。
  解説を読んでから演習に取り組み、模範解答と見比べてください。出席確認クイズは、講義サイトの出席確認ページで
  実際に解答したうえで、解説を読んで理解を確かめてください。
\end{tcolorbox}

\vspace{3mm}
{\sffamily\large\bfseries\color{mone}目次}\par\vspace{0.5mm}
\makeatletter
\renewcommand\contentsname{}
\makeatother
{\small\tableofcontents}
